\section{Beurteilung}
\label{typ:beurteilung}
Bei der Nutzung des Dokumenttyps Beurteilung werden einige Makros geladen, so dass Beurteilungen von
Lehramtsanwärterinnen und Lehramtsanwärtern bzw. Lehrkräften in Ausbildung einfach gesetzt werden
können. U.\,a. wird dazu die Titelseite entsprechend gestaltet und Umgebungen für die
Handlungsfelder vorformatiert.

\subsection{Paketoptionen}
Die folgenden Optionen sollten in der Präambel dem Schule-Paket mitgegeben werden.
\begin{options}
 \opt{beurteilung}
   Gibt den Namen des Beurteilenden an, z.\,B. \texttt{beurteilung=\{Marine Musterfrau\}}.

 \opt{ref}
   Gibt den Namen des Beurteilten an, z.\,B. \texttt{ref=\{Sabine Musterref\}}.

 \opt{zeitraum}
	 Gibt den Zeitraum des Gutachtens an (nicht nur die einzelnen gesehenen Unterrichte), z.\,B.
\texttt{zeitraum=\{02.02.19--03.05.19\}}.

 \opt{schulname}
	Gibt den Namen der Schule an.

 \opt{schullogo}
  Lädt das Schullogo auf die Titelseite rechts oben. Das Schullogo sollte \LaTeX als
\texttt{logo.pdf} vorliegen.

 \opt{schulanschrift}
  Wenn ein Schullogo gesetzt wird, kann zusätzlich der Name der Schule,

 \opt{schulstr}
  die Straße der Schule und

 \opt{schulort}
  der Ort der Schule ausführlich unter das Logo geschrieben werden.

 \opt{lehramt}
  Gibt an, um welches Lehramt es sich handelt, z.\,B. GyGe.
\end{options}

 \subsection{Befehle}
 \begin{commands}
     \command{setzeGrundlagen}[\{Termine\}]
				 Über den Mechanismus können auf der Titelseite Einträge in die Tabelle zum hospitierten
bzw. gegebenen Unterricht eingetragen werden. Dabei werden drei Spalten erwartet. Z.\,B.:

			\texttt{\textbackslash setzeGrundlagen\{02.04.--01.05.2019 \& 9b \& Themen\}}
 \end{commands}
 
\subsection{Umgebungen}
 \begin{commands}
	 \command{begin}[\{handlungsfeld1\}]
	 \command{begin}[\{handlungsfeld2\}]
	 \command{begin}[\{handlungsfeld3\}]
	 \command{begin}[\{handlungsfeld4\}]
	 \command{begin}[\{handlungsfeld5\}]
	 \command{begin}[\{handlungsfeld6\}]

	In den Umgebungen \texttt{handlungsfeldX} können die Beurteilungen entsprechend zu den
Handlungsfeldern geschrieben werden. Jedes Handlungsfeld wird dazu in einer \texttt{tcolorbox}
gesetzt. Die zugehörigen Kompetenzen und Titel werden passend zugeordnet. Optional können der
	 Umgebung weitere Optionen für die tcolorbox mitgegeben werden:

	 \texttt{\textbackslash begin\{handlungsfeld1\}[colback=red]}.
 \end{commands}
